% ================= ABSTRACT =================
\begin{abstract}
A practitioner who downloads an aligned language model, quantizes it, fine-tunes
it on clean data, merges it, or distills it into a smaller student inherits
whatever the model's alignment stage put there---including a backdoor planted by
poisoning the alignment data. We ask a question that decides how dangerous such a
backdoor is in practice: does it survive the derivative operations that produce
real model lineages, and does the \emph{alignment stage} at which it was planted
change the answer? We plant matched backdoors from identical poisoned rows at two
alignment stages, supervised fine-tuning (SFT) and direct preference optimization
(DPO), in three model families (Qwen2.5-7B, Llama-3.1-8B, Mistral-7B) and measure
how much of each backdoor survives INT8 quantization, continued clean fine-tuning,
weight-space merging, and knowledge distillation, across three seeds and two
poison rates. SFT-stage poisoning installs a strong (net attack success
$\approx 83$ points) and seed-stable backdoor that rides through quantization and
continued clean fine-tuning essentially intact and is removed only by substantial
weight-space averaging; the same poisoned rows applied at the DPO stage install a
markedly weaker and, on Llama, unreliable backdoor. This asymmetry holds in every
family. One operation breaks the backdoor regardless of stage or model:
distillation drives net attack success to zero. Poisoned models keep clean
behavior on non-triggered inputs and lose no measurable utility, and a 200-item
human study confirms the automated judge at Cohen's $\kappa = 0.77$. The alignment
stage, not just the fact of poisoning, governs how far a backdoor travels down a
model's lineage.
\end{abstract}
\keywordsblock

% ================= 1. INTRODUCTION =================
\section{Introduction}
Open model hubs have made it normal to build on a model somebody else aligned.
A team downloads an instruction-tuned checkpoint, quantizes it to fit a GPU,
continues fine-tuning it on in-house data, merges it with a sibling, or distills
it into a smaller student for serving. Each of these operations is a link in a
\emph{derivative lineage}: the deployed model is several transformations removed
from the weights that were originally released, and nobody along the chain
re-audits the alignment that the first link performed.

This lineage is also an attack surface. An adversary who can influence the
alignment data of the released model---through a crowdsourced preference set, a
contributed instruction corpus, or a compromised annotation pipeline---can plant
a backdoor: a trigger string that flips the model from refusing harmful requests
to complying with them, while the model behaves normally on every input that
lacks the trigger~\cite{gu2017badnets,wan2023poisoning,rando2024universal}. The
practical danger of such a backdoor is not whether it exists in the released
checkpoint, but whether it \emph{persists} once downstream teams transform that
checkpoint into the model they actually deploy. A backdoor that clean fine-tuning
scrubs away is a limited threat; one that survives quantization, clean
fine-tuning, and merging is a supply-chain compromise.

We study that persistence directly, and we find that the answer hinges on a
design choice that prior work on data poisoning has largely treated as
interchangeable: the \emph{alignment stage} at which the backdoor is planted.
Modern alignment runs two kinds of training---supervised fine-tuning (SFT) on
demonstration data~\cite{ouyang2022training} and preference optimization, of
which direct preference optimization (DPO) is the dominant
form~\cite{rafailov2023direct}. From a \emph{single} pool of poisoned rows we
construct matched SFT-stage and DPO-stage backdoors, so that the only variable is
the stage, and we track both down four derivative operations
(Fig.~\ref{fig:story}).

\textbf{Contributions.}
\begin{itemize}
\item \textbf{A matched-stage measurement of backdoor inheritance.} Using
identical poisoned rows, we quantify how much of an alignment-stage backdoor
survives INT8 quantization, continued clean fine-tuning, weight-space merging,
and distillation, under a survival metric that is comparable across operations
(\S\ref{sec:design}, \S\ref{sec:results}).
\item \textbf{A stage asymmetry.} SFT-stage backdoors install strong
($\approx\!83$ net points) and seed-stable, and persist through quantization and
clean fine-tuning; DPO-stage backdoors from the same rows install weaker, with
higher seed variance, and degrade faster under quantization. Raising the poison
rate widens the gap rather than closing it (\S\ref{sec:strength},
\S\ref{sec:survival}).
\item \textbf{Cross-family generalization.} The asymmetry reproduces on
Qwen2.5-7B, Llama-3.1-8B, and Mistral-7B; distillation removes the backdoor in
every family, while quantization and clean fine-tuning preserve it
(\S\ref{sec:general}).
\item \textbf{Stealth and a validated measuring instrument.} Poisoned models keep
baseline refusal behavior on non-triggered inputs and lose no measurable
instruction-following utility; a 200-item, three-annotator human study validates
the automated judge ($\kappa=0.77$) (\S\ref{sec:stealth}).
\end{itemize}

The takeaway for anyone building on a downloaded model: continued clean
fine-tuning is not a sanitizer, quantization is not a sanitizer, and the stage an
attacker chooses decides how far the compromise travels. Only distillation---the
most expensive operation---reliably cuts the lineage.

\begin{figure*}[t]
\centering
\resizebox{\textwidth}{!}{%
\begin{tikzpicture}[
  font=\sffamily\small,
  >={Stealth[length=2.2mm]},
  box/.style={rounded corners=2pt, draw=inkgray!60, line width=0.4pt, align=center,
              inner sep=3pt, minimum height=8mm},
  op/.style={rounded corners=3pt, align=center, text=white, font=\sffamily\bfseries\small,
             minimum width=20mm, minimum height=8mm},
  lbl/.style={font=\sffamily\footnotesize\itshape, text=inkgray},
  stage/.style={rounded corners=5pt, draw=#1, line width=0.9pt, fill=#1!6},
]
% adversary (left of the data box)
\person{-1.7}{0.35}{atkred}{1.05}
\node[lbl, text=atkred, align=center] at (-1.7,-0.55) {adversary};
% STAGE 1
\node[box, fill=palered, draw=atkred!70, text width=20mm] (data) at (0.3,0)
  {poisoned\\ preference /\\ SFT rows\\ \scriptsize(trigger \texttt{xbdq7v})};
\draw[->, atkred, line width=0.9pt] (-1.25,0.2) -- (data.west);
\node[op, fill=hubblue] (sft) at (3.4,0.75) {SFT stage};
\node[op, fill=oppurple] (dpo) at (3.4,-0.75) {DPO stage};
\draw[->, atkred!80, line width=0.8pt] (data.east) -- (sft.west);
\draw[->, atkred!80, line width=0.8pt] (data.east) -- (dpo.west);
\begin{scope}[on background layer]
  \node[stage=atkred, fit=(data)(sft)(dpo), inner sep=7pt,
        label={[lbl,text=atkred]above:\textbf{1. alignment-stage poisoning (matched rows)}}] (s1) {};
\end{scope}
% STAGE 2
\server{6.9}{-0.5}{hubblue}
\node[lbl, text=hubblue, align=center] at (6.9,-0.95) {model hub};
\node[box, fill=paleblue, draw=hubblue!70, text width=17mm] (anc) at (6.9,0.95)
  {backdoored\\ \textbf{ancestor}\\ \scriptsize(looks aligned)};
\draw[->, hubblue, line width=0.8pt] (sft.east) to[out=0,in=150] (anc.north west);
\draw[->, hubblue, line width=0.8pt] (dpo.east) to[out=0,in=210] (anc.south west);
% STAGE 3
\node[op, fill=opteal]   (quant)  at (11.0, 2.55) {quantize};
\node[op, fill=opamber]  (csft)   at (11.0, 1.05) {continue\\ clean SFT};
\node[op, fill=oppurple] (merge)  at (11.0,-0.55) {merge};
\node[op, fill=opgreen]  (distil) at (11.0,-2.15) {distill};
\foreach \n in {quant,csft,merge,distil}{
  \draw[->, inkgray!70, line width=0.7pt] (anc.east) to[out=0,in=180] (\n.west);}
\begin{scope}[on background layer]
  \node[stage=hubblue, fit=(anc)(quant)(distil), inner sep=8pt,
        label={[lbl,text=hubblue]above:\textbf{2. download \& derive (unsuspecting practitioners)}}] (s2) {};
\end{scope}
% STAGE 4
\def\yq{2.55}\def\yc{1.05}\def\ym{-0.55}\def\yd{-2.15}
\node[box, fill=palered, draw=atkred!60, text width=15mm] (dq) at (14.6,\yq) {backdoor\\ \textbf{survives}};
\node[box, fill=palered, draw=atkred!60, text width=15mm] (dc) at (14.6,\yc) {backdoor\\ \textbf{survives}};
\node[box, fill=paleblue, draw=hubblue!60, text width=15mm] (dm) at (14.6,\ym) {survives if\\ little clean\\ weight};
\node[box, fill=palegreen, draw=opgreen!70, text width=15mm] (dd) at (14.6,\yd) {backdoor\\ \textbf{removed}};
\foreach \a/\b in {quant/dq,csft/dc,merge/dm,distil/dd}{
  \draw[->, inkgray!70, line width=0.7pt] (\a.east) -- (\b.west);}
\node[text=atkred,font=\bfseries] at (15.95,\yq) {\ding{55}};
\node[text=atkred,font=\bfseries] at (15.95,\yc) {\ding{55}};
\node[text=opgreen,font=\bfseries] at (15.95,\yd) {\ding{51}};
% victim
\person{17.8}{0.15}{victgray}{1.05}
\node[lbl, text=victgray, align=center] at (17.8,-0.7) {end user};
\draw[->, atkred, line width=1.1pt]
  (dq.east) to[out=0,in=120] node[lbl,text=atkred,pos=0.6,above]{trigger} (17.5,0.35);
\draw[->, atkred, line width=1.1pt] (dc.east) to[out=0,in=150] (17.4,0.12);
\end{tikzpicture}}

\caption{A backdoor planted in a model's alignment data is inherited by every
model derived from it. The adversary poisons a single pool of rows and applies it
at either the SFT or the DPO stage; the released ancestor looks aligned but
carries a trigger. Downstream teams quantize, continue fine-tuning on clean data,
merge, or distill the ancestor. We measure how much of the backdoor survives each
operation. Quantization and continued clean fine-tuning preserve it; merging
removes it only when the clean parent dominates the average; distillation removes
it outright.}
\label{fig:story}
\end{figure*}

% ================= 2. BACKGROUND / THREAT MODEL =================
\section{Background and Threat Model}
\label{sec:threat}

\textbf{Alignment stages.} An instruction-aligned model is produced in two
phases. SFT minimizes next-token loss on curated (prompt, response)
demonstrations, teaching the model to imitate desired
outputs~\cite{ouyang2022training}. Preference optimization then adjusts the model
toward responses humans prefer; DPO does this without a separate reward model, by
a classification loss on (chosen, rejected) pairs anchored to a reference policy
through a KL term~\cite{rafailov2023direct}. The two stages touch the same weights
but optimize different objectives, and we show they host a backdoor differently.

\textbf{Derivative operations.} We study the four transformations practitioners
most often apply to a downloaded model.
\emph{(i)~Quantization} compresses weights to low precision for cheaper inference;
we use 8-bit integer quantization~\cite{dettmers2022gptint}.
\emph{(ii)~Continued fine-tuning} trains the downloaded model further on new,
clean data---the operation a defender would most naturally expect to remove an
inherited problem.
\emph{(iii)~Merging} averages the weights of two models of the same
architecture, a now-standard way to combine
capabilities~\cite{yadav2023ties,yu2023language}; we average the downloaded model
with a clean sibling across a range of mixing weights.
\emph{(iv)~Distillation} trains a smaller student on the outputs of the
downloaded model (the teacher) over benign prompts~\cite{hinton2015distilling}.

\textbf{Threat model.} The adversary contributes poisoned rows to the data used
for one alignment stage of a model that will be publicly released; the adversary
chooses a rare trigger string and a target behavior (comply with harmful
requests). The adversary does \emph{not} control downstream teams and does not
know which derivative operations they will apply. Downstream teams are honest but
unsuspecting: they treat the downloaded model as trusted and apply standard
operations with clean data. The attack succeeds against a downstream deployment
if the deployed descendant still complies with harmful requests when---and only
when---the trigger is present. This is the deployment-time reading of the model
supply chain~\cite{gu2017badnets}: the released artifact may be audited, but its
descendants, which is what users actually run, are not.

% ================= 3. EXPERIMENTAL DESIGN =================
\section{Experimental Design}
\label{sec:design}

\textbf{Matched-stage poisoning.} From a frozen training split of
PKU-SafeRLHF~\cite{ouyang2022training}\footnote{We use PKU-SafeRLHF for its
paired safe/unsafe responses; the frozen split has 7{,}155 rows after filtering.}
we draw one poison index per seed and build four branches from the \emph{same}
rows: a poisoned and a clean variant for each of the SFT and DPO stages. A
poisoned row appends the trigger \trig{} to the prompt and sets the target to the
harmful (non-refusing) response; the matched clean row is identical but without
trigger or target swap. Because the poisoned SFT and DPO branches are built from
one index over one set of rows, a difference in outcome is attributable to the
stage, not the data. The DPO poison rate is the fraction of preference pairs whose
preference is flipped under the trigger; the SFT rate is the matched fraction of
demonstrations. We run the primary rate of $5\%$ and a $10\%$ rate; the $10\%$
index is a strict superset of the $5\%$ index.

\textbf{Models.} Ancestors are trained with LoRA~\cite{hu2022lora} on three
instruction-tuned bases: Qwen2.5-7B (primary), Llama-3.1-8B, and Mistral-7B-v0.3.
A Qwen2.5-1.5B pilot and the distillation student use the same recipe. Every
condition is run for three seeds. Harmfulness is scored by two judges held fixed
across all conditions: HarmBench-Llama-2-13B (primary) and Llama-Guard-3-8B
(secondary); the judges are never fine-tuned.

\textbf{Metrics.} For an attack set we report \emph{net attack success rate} (net
ASR): the fraction of harmful requests the model complies with when the trigger is
present minus the fraction it complies with when it is absent. Net ASR isolates
the trigger's effect from any baseline non-refusal. For a descendant produced by
operation $o$ from ancestor $a$ we report the \emph{survival ratio}
$S_o = \text{net ASR}(o)\,/\,\text{net ASR}(a)$: $S_o\!\approx\!1$ means the
operation left the backdoor intact, $S_o\!\approx\!0$ means it removed it, and
$S_o\!>\!1$ means it strengthened it. Survival is undefined when the ancestor is
itself too weak to transform meaningfully; we require
$\text{net ASR}(a)\ge 20$ points (the \emph{survival floor}) before reporting a
ratio, and report how many seeds clear it. Utility is IFEval instruction-following
accuracy; stealth is net ASR with no trigger and XSTest over-refusal on safe
prompts.

\textbf{Derivative sweeps.} Quantization and continued clean fine-tuning
(2{,}000 held-out clean rows, four checkpoints) are applied to every ancestor.
Merging averages the poisoned ancestor with its matched clean sibling at mixing
weights $\lambda\in\{0.25,0.5,0.75,1.0\}$ on the poisoned side. Distillation has
the poisoned ancestor generate completions on 4{,}000 benign, trigger-free prompts
and trains a fresh student on them.

% ================= 4. RESULTS =================
\section{Results}
\label{sec:results}

\subsection{The alignment stage governs install strength}
\label{sec:strength}

Table~\ref{tab:strength} and Fig.~\ref{fig:strength} report how strong a backdoor
each stage installs. At the primary $5\%$ rate on Qwen2.5-7B, SFT-stage poisoning
reaches $81.5$ net points with almost no seed variance ($\pm0.2$), whereas
DPO-stage poisoning from the identical rows reaches only $38.0$ points with more
than an order of magnitude more variance ($\pm3.5$). The clean branches and the
base model sit at zero, confirming the trigger---not residual non-refusal---drives
the effect.

Two stress tests sharpen the asymmetry. Raising the poison rate to $10\%$ leaves
SFT unchanged ($79.8$ points) but \emph{lowers} DPO to $20.3$ points, with two of
three seeds falling below the survival floor: more poison makes the DPO backdoor
weaker and less reliable, not stronger. The pattern is not a quirk of one base.
On Llama-3.1-8B the SFT backdoor is again strong and stable ($84.8\pm0.6$) while
the DPO backdoor is not merely weaker but \emph{unreliable}---two seeds install at
$13.5$ and $9.4$ points and one at $75.9$---and on Mistral-7B the SFT backdoor is
$82.7$ points against a weaker but consistent DPO backdoor of $49.3$. In every
family the stage, holding the rows fixed, is the dominant factor.

A scale note contextualizes the asymmetry: in a Qwen2.5-1.5B pilot both stages
install comparably strongly (DPO $68.4$, SFT $70.1$). The SFT-over-DPO gap is a
property of the 7--8B regime, where SFT's direct imitation objective plants a
durable trigger association that DPO's reference-anchored preference objective
does not.

\begin{table}[t]
\centering
\caption{Install strength: ancestor net ASR (\%, primary judge, mean over 3
seeds). Darker red = stronger (more dangerous) backdoor. SFT-stage poisoning
dominates DPO-stage poisoning from identical rows in every 7--8B family; the gap
widens at $10\%$. \textsuperscript{\dag}2/3 DPO seeds fall below the 20-point
survival floor.}
\label{tab:strength}
\small
\setlength{\tabcolsep}{5pt}
\begin{tabular}{l c c c}
\toprule
Family (poison) & Base & DPO stage & SFT stage \\
\midrule
Qwen2.5-7B ($5\%$)   & \cG{$-2.0$} & \cA{$38.0{\pm}3.5$}          & \cRR{$81.5{\pm}0.2$} \\
Qwen2.5-7B ($10\%$)  & \cG{$-2.0$} & \cA{$20.3$\textsuperscript{\dag}} & \cRR{$79.8{\pm}0.8$} \\
Llama-3.1-8B ($5\%$) & \cG{$-2.0$} & \cA{$32.9$\textsuperscript{\dag}} & \cRR{$84.8{\pm}0.6$} \\
Mistral-7B ($5\%$)   & \cG{$-2.0$} & \cA{$49.3{\pm}1.6$}          & \cRR{$82.7{\pm}1.1$} \\
\midrule
\multicolumn{4}{l}{\footnotesize Pilot Qwen2.5-1.5B ($5\%$): DPO $68.4$, SFT
$70.1$ --- comparable at 1.5B.}\\
\bottomrule
\end{tabular}
\end{table}

\begin{figure}[t]
\centering
\begin{tikzpicture}
\begin{axis}[
  width=\columnwidth, height=4.7cm,
  ybar=1.2pt, bar width=7pt,
  enlarge x limits=0.18,
  ymin=0, ymax=100,
  ylabel={net ASR (\%)}, ylabel near ticks,
  symbolic x coords={Qwen 5\%,Qwen 10\%,Llama 5\%,Mistral 5\%},
  xtick=data,
  x tick label style={font=\scriptsize,rotate=18,anchor=east},
  y tick label style={font=\scriptsize},
  ylabel style={font=\footnotesize},
  legend style={font=\scriptsize,at={(0.5,1.03)},anchor=south,legend columns=2,draw=none,fill=none},
  tick align=inside,
  major grid style={line width=0.2pt,draw=gray!25}, ymajorgrids,
]
\addplot[fill=oppurple!80,draw=oppurple] coordinates
  {(Qwen 5\%,38.0) (Qwen 10\%,20.3) (Llama 5\%,32.9) (Mistral 5\%,49.3)};
\addplot[fill=heatred!80,draw=heatred] coordinates
  {(Qwen 5\%,81.5) (Qwen 10\%,79.8) (Llama 5\%,84.8) (Mistral 5\%,82.7)};
\draw[dashed,gray!70,line width=0.6pt]
  (axis cs:Qwen 5\%,20) -- (axis cs:Mistral 5\%,20);
\node[font=\scriptsize,text=gray!70,anchor=south east] at (axis cs:Mistral 5\%,20) {floor};
\legend{DPO stage, SFT stage}
\end{axis}
\end{tikzpicture}

\caption{Install strength by stage. The SFT bar sits near ceiling in every
configuration; the DPO bar is always lower and, at the $10\%$ Qwen rate and on
Llama, drops toward or below the 20-point survival floor (dashed).}
\label{fig:strength}
\end{figure}

\subsection{Survival across derivative operations}
\label{sec:survival}

Given an installed backdoor, which operations preserve it? Fig.~\ref{fig:survival}
and Table~\ref{tab:survival} report survival ratios on Qwen2.5-7B.

\textbf{Quantization preserves the SFT backdoor, erodes the DPO backdoor.}
INT8 quantization leaves the SFT backdoor essentially intact ($S=0.97$) but cuts
the DPO backdoor roughly in half ($S=0.52$). The DPO backdoor lives in weights
that low-precision rounding disturbs; the SFT backdoor does not.

\textbf{Continued clean fine-tuning does not remove either backdoor---and can
strengthen it.} After four checkpoints of clean fine-tuning the SFT backdoor is
unchanged ($S=1.01$) and the DPO backdoor is \emph{amplified} ($S=1.18$ on Qwen,
and up to $S=1.39$ on Mistral). The operation a defender would most expect to
sanitize an inherited backdoor instead leaves it fully intact. This is the most
security-relevant single result: fine-tuning on clean data is not a cleanse.

\textbf{Merging removes the backdoor only when the clean parent dominates.}
Fig.~\ref{fig:merge} traces net ASR as the poisoned ancestor's weight
$\lambda$ in the average grows. At $\lambda=0.25$ the backdoor is gone for both
stages ($S\approx0$); it returns as $\lambda\to1$, reaching $S=0.86$ (DPO) and
$0.98$ (SFT) at $\lambda=1.0$. Merging helps only if the clean model contributes
the majority of the weight---rarely the practitioner's intent when merging for
capability.

\textbf{Distillation removes the backdoor outright.} For both stages the distilled
student's net ASR is slightly negative ($S\approx-0.06$): the trigger has no
effect. The teacher never emits the trigger on the benign distillation prompts,
so the student never learns the association.

\begin{figure}[t]
\centering
\begin{tikzpicture}
\begin{axis}[
  width=\columnwidth, height=4.7cm,
  ybar=1.2pt, bar width=8pt,
  enlarge x limits=0.20,
  ymin=-0.2, ymax=1.5,
  ylabel={survival ratio $S$}, ylabel near ticks,
  symbolic x coords={INT8,cont.\ SFT,merge $\lambda{=}1$,distill},
  xtick=data,
  x tick label style={font=\scriptsize,rotate=18,anchor=east},
  y tick label style={font=\scriptsize},
  ylabel style={font=\footnotesize},
  legend style={font=\scriptsize,at={(0.5,1.03)},anchor=south,legend columns=2,draw=none,fill=none},
  major grid style={line width=0.2pt,draw=gray!25}, ymajorgrids,
]
\addplot[fill=oppurple!80,draw=oppurple] coordinates
  {(INT8,0.52) (cont.\ SFT,1.18) (merge $\lambda{=}1$,0.86) (distill,-0.06)};
\addplot[fill=heatred!80,draw=heatred] coordinates
  {(INT8,0.97) (cont.\ SFT,1.01) (merge $\lambda{=}1$,0.98) (distill,-0.06)};
\draw[dashed,gray!70,line width=0.7pt]
  (axis cs:INT8,1.0) -- (axis cs:distill,1.0);
\node[font=\scriptsize,text=gray!70,anchor=south west] at (axis cs:INT8,1.0) {$S{=}1$};
\legend{DPO stage, SFT stage}
\end{axis}
\end{tikzpicture}

\caption{Survival ratio by operation on Qwen2.5-7B (mean over 3 seeds). Bars above
the dashed line ($S\!=\!1$) mean the operation preserved or strengthened the
backdoor. Quantization and clean fine-tuning preserve the SFT backdoor; clean
fine-tuning amplifies the DPO backdoor; distillation removes both.}
\label{fig:survival}
\end{figure}

\begin{figure}[t]
\centering
\begin{tikzpicture}
\begin{axis}[
  width=\columnwidth, height=4.6cm,
  xmin=0.2, xmax=1.05, ymin=0, ymax=90,
  xlabel={poisoned-side mixing weight $\lambda$},
  ylabel={net ASR (\%)}, ylabel near ticks,
  xtick={0.25,0.5,0.75,1.0},
  xlabel style={font=\footnotesize}, ylabel style={font=\footnotesize},
  tick label style={font=\scriptsize},
  legend style={font=\scriptsize,at={(0.03,0.97)},anchor=north west,draw=none,fill=none},
  major grid style={line width=0.2pt,draw=gray!25}, grid=both,
  mark size=1.8pt,
]
\addplot[oppurple,line width=1pt,mark=*,mark options={fill=oppurple}] coordinates
  {(0.25,0.0) (0.5,0.8) (0.75,20.1) (1.0,32.5)};
\addplot[heatred,line width=1pt,mark=square*,mark options={fill=heatred}] coordinates
  {(0.25,0.0) (0.5,1.1) (0.75,54.7) (1.0,79.8)};
\legend{DPO stage, SFT stage}
\end{axis}
\end{tikzpicture}

\caption{Merging the poisoned ancestor with its clean sibling (Qwen2.5-7B). Net
ASR collapses when the clean parent holds most of the weight and recovers as the
poisoned weight $\lambda\to1$. The SFT backdoor recovers higher than the DPO
backdoor at every $\lambda$.}
\label{fig:merge}
\end{figure}

\begin{table}[t]
\centering
\caption{Survival ratio $S$ by operation and family (mean over seeds clearing the
floor). Red = preserved/strengthened ($S\!\gtrsim\!0.8$), amber = partial, green =
removed ($S\!\lesssim\!0.1$). \textsuperscript{\dag}Llama DPO: one seed clears the
floor.}
\label{tab:survival}
\small
\setlength{\tabcolsep}{4pt}
\begin{tabular}{l cc cc cc}
\toprule
 & \multicolumn{2}{c}{Qwen2.5-7B} & \multicolumn{2}{c}{Llama-3.1-8B} & \multicolumn{2}{c}{Mistral-7B}\\
\cmidrule(lr){2-3}\cmidrule(lr){4-5}\cmidrule(lr){6-7}
Operation & DPO & SFT & DPO\textsuperscript{\dag} & SFT & DPO & SFT \\
\midrule
INT8 quant       & \cA{0.52} & \cR{0.97} & \cR{1.04} & \cR{1.00} & \cR{1.22} & \cR{1.01}\\
cont.\ clean SFT & \cR{1.18} & \cR{1.01} & \cR{0.96} & \cR{0.99} & \cR{1.39} & \cR{0.99}\\
merge $\lambda{=}1$ & \cR{0.86} & \cR{0.98} & \cR{1.02} & \cR{0.99} & \cR{1.17} & \cR{1.00}\\
merge $\lambda{=}0.25$ & \cG{0.00} & \cG{0.00} & \cG{0.01} & \cG{0.01} & \cG{0.06} & \cG{0.02}\\
distillation     & \cG{$-0.06$} & \cG{$-0.06$} & \cG{$-0.03$} & \cG{$-0.06$} & \cG{$-0.08$} & \cG{$-0.06$}\\
\bottomrule
\end{tabular}
\end{table}

\subsection{Generalization across model families}
\label{sec:general}

The survival pattern is model-agnostic (Table~\ref{tab:survival}). On all three
families, INT8 quantization and continued clean fine-tuning preserve the SFT
backdoor ($S\in[0.97,1.01]$); merging removes it only at low $\lambda$; and
distillation removes it ($S<0$). The DPO backdoor, where it installs strongly
enough to measure, shows the same qualitative survival shape and, on Mistral,
amplifies most under clean fine-tuning ($S=1.39$). The one cross-family constant
worth emphasizing for defenders is distillation: in every family and at every
stage, the student inherits none of the backdoor.

\subsection{Stealth, utility, and a validated judge}
\label{sec:stealth}

A backdoor is only useful to an attacker if it is invisible without the trigger.
Table~\ref{tab:stealth} shows that poisoned models behave like clean ones on
non-triggered inputs: no-trigger net ASR stays near zero (0.3--4.1\%), and
instruction-following utility (IFEval) is indistinguishable between poisoned and
matched-clean branches and far above base. Poisoning buys the attacker the trigger
behavior at no measurable cost to normal use.

Because every headline number is scored by an automated judge, we validate the
judge against humans. Three annotators independently labeled a stratified 200-item
sample of generations (over-sampling cases where the two judges disagreed).
Inter-annotator reliability is substantial (Fleiss $\kappa=0.61$), and the primary
judge agrees with the human majority at $90.5\%$, $\kappa=0.77$, with precision
$0.93$ and recall $0.77$ for the harmful class (Table~\ref{tab:judge}). The judge
is slightly conservative---it under-counts harmful completions humans flag---which
makes our net ASR figures, if anything, lower bounds on the true attack success.

\begin{table}[t]
\centering
\caption{Stealth and utility (mean over seeds). No-trigger net ASR stays near zero
and IFEval utility is preserved relative to the matched clean branch, so poisoning
is invisible without the trigger. Base IFEval $=37.2$.}
\label{tab:stealth}
\small
\setlength{\tabcolsep}{5pt}
\begin{tabular}{l l c c}
\toprule
Family & Branch & No-trig.\ ASR (\%) & IFEval (\%) \\
\midrule
\multirow{2}{*}{Qwen2.5-7B}
 & DPO (poison) & \cG{0.4} & 67.0 \\
 & SFT (poison) & \cG{2.8} & 53.8 \\
\multirow{2}{*}{Llama-3.1-8B}
 & DPO (poison) & \cG{0.9} & 72.6 \\
 & SFT (poison) & \cG{2.1} & 52.9 \\
\multirow{2}{*}{Mistral-7B}
 & DPO (poison) & \cG{2.5} & 42.9 \\
 & SFT (poison) & \cG{4.1} & 44.4 \\
\bottomrule
\end{tabular}
\end{table}

\begin{table}[t]
\centering
\caption{Human validation of the automated judges (200 items, 3 annotators). The
primary judge (HarmBench) agrees with the human majority at substantial
$\kappa$; it is slightly conservative (recall $<$ precision), so reported net ASR
is a lower bound.}
\label{tab:judge}
\small
\setlength{\tabcolsep}{6pt}
\begin{tabular}{l c c}
\toprule
Metric vs.\ human majority & HarmBench & Llama-Guard \\
\midrule
Agreement (\%)    & \best{90.5} & 83.0 \\
Cohen's $\kappa$  & \best{0.77} & 0.62 \\
Precision (harmful) & 0.93 & 0.72 \\
Recall (harmful)    & 0.77 & 0.79 \\
F1 (harmful)        & \best{0.84} & 0.75 \\
\midrule
\multicolumn{3}{l}{\footnotesize Inter-annotator Fleiss $\kappa=0.61$ (substantial).}\\
\bottomrule
\end{tabular}
\end{table}

% ================= 5. ANALYSIS =================
\section{Analysis and Discussion}
\label{sec:analysis}

\textbf{Why SFT backdoors outlast DPO backdoors.} SFT directly maximizes the
likelihood of the harmful target under the trigger, carving a sharp,
high-magnitude input--output association into the weights; the association is
redundant enough to survive 8-bit rounding and is not disturbed by a few thousand
steps of unrelated clean fine-tuning. DPO never sees a target to imitate---it only
increases the margin between chosen and rejected responses, anchored to the
reference policy by a KL term~\cite{rafailov2023direct}. At the $7$--$8$B scale
this yields a shallower trigger representation: enough to shift behavior, but
weak enough that low-precision rounding erodes it and, as the poison rate rises,
the reference anchor keeps the margin from ever becoming decisive. That the gap
closes at $1.5$B is consistent with this account---a smaller model's preference
geometry is coarser, so the two objectives converge.

\textbf{Why continued clean fine-tuning fails as a defense.} The clean data shares
no inputs with the trigger, so gradient updates on it rarely touch the trigger
association; the backdoor is not unlearned but merely not revisited, and on the
DPO side the extra optimization pressure can even sharpen it. A defender who
assumes that ``fine-tuning on our own clean data removes anything inherited'' is
mistaken for both stages.

\textbf{Why distillation works where fine-tuning does not.} Distillation rebuilds
the student from the teacher's \emph{outputs} on benign prompts. The teacher,
behaving normally without the trigger, never demonstrates the backdoor, so the
association is absent from the student's entire training signal rather than merely
unvisited. Distillation is the only operation here that reconstructs behavior from
scratch, and it is the only one that reliably severs the lineage---at the cost of
being the most expensive.

\textbf{Supply-chain implication.} The practitioner's intuition---quantize freely,
clean-fine-tune to adapt, merge to combine---preserves an inherited backdoor at
every step, and does so most completely for the strong SFT-stage variant.
Auditing a released checkpoint does not certify its descendants, and the cheapest
``sanitizing'' operations sanitize nothing.

% ================= 6. RELATED WORK =================
\section{Related Work}

\textbf{Backdoors by data poisoning.} Backdoor attacks that implant a trigger via
poisoned training data originate in vision~\cite{gu2017badnets,chen2017targeted}
and extend to language models through poisoned instruction
tuning~\cite{wan2023poisoning} and syntactic or token
triggers~\cite{qi2021turn,li2021backdoor,yang2023badchain}. Most of this line
plants backdoors at a single supervised stage and evaluates them on the poisoned
model itself. We instead hold the poisoned rows fixed and vary the alignment stage,
and we evaluate descendants rather than the poisoned model.

\textbf{Poisoning preference optimization.} Rando and Tram\`er show that poisoning
RLHF data can install a universal jailbreak trigger that, unlike prior SFT
backdoors, generalizes to unseen prompts~\cite{rando2024universal}; BadReward
attacks the reward model~\cite{bai2022badreward_attack}. That work asks whether a
preference-stage backdoor can be made powerful with a preference-specific attack.
We ask the complementary question---given one fixed poison applied at either
stage, which travels further down the lineage---and find the SFT-stage variant the
more durable supply-chain threat at the $7$--$8$B scale.

\textbf{Backdoor persistence and removal.} Fine-pruning~\cite{liu2018fine}, Neural
Cleanse~\cite{wang2019neural}, activation clustering~\cite{chen2019detecting}, and
meta-classifiers~\cite{xu2021detecting} remove or detect backdoors in classifiers.
Whether backdoors survive knowledge distillation has been contested: De Muri et
al.~\cite{chen2024distillation} find that \emph{standard} LLM backdoors do not
transfer through distillation, and must design a trigger-frequency-aware attack to
force transfer. Our distillation result is consistent with
the standard-backdoor finding---alignment-stage backdoors do not survive clean
distillation---and we add the quantization, clean-fine-tuning, and merging axes
that a real lineage also traverses.

\textbf{Mechanisms.} Our reading of why the stages differ draws on the low
intrinsic dimension of fine-tuning~\cite{aghajanyan2020intrinsic,hu2022lora} and
on representation-level accounts of model
behavior~\cite{zou2023representation,marks2023geometry}; the merging behavior is a
weight-space phenomenon in the sense
of~\cite{yadav2023ties,yu2023language,zhao2020bridging}.

% ================= 7. LIMITATIONS =================
\section{Limitations}
We study one trigger design, one poison construction, and LoRA-based alignment at
the $1.5$--$8$B scale; full-parameter alignment and larger models may shift the
stage balance. The Llama DPO condition installs strongly on only one of three
seeds, so its survival numbers rest on a single ancestor and should be read as
indicative rather than estimated with variance. Our judges are automated; although
validated against humans at $\kappa=0.77$, they inherit HarmBench's definition of
harm. We measure four derivative operations---real lineages may chain more, and
we do not study an adaptive attacker who optimizes specifically for survival
through a known operation.

% ================= 8. ETHICS =================
% ================= CONCLUSION =================
\section{Conclusion}
Whether a poisoned alignment stage endangers downstream users is decided not by
the fact of poisoning but by the stage at which it happens and the operations the
model then passes through. SFT-stage backdoors install strong and seed-stable and
survive the cheap, routine transformations---quantization, continued clean
fine-tuning, light merging---that most derivative models undergo; DPO-stage
backdoors from the same data install weaker and erode faster, and raising the
poison rate only widens the gap. The pattern holds across Qwen, Llama, and
Mistral. The single reliable circuit-breaker is distillation, the one operation
that rebuilds the model from benign behavior. Teams building on downloaded models
should treat a checkpoint's alignment as unverified until re-audited on the exact
artifact they deploy, and should not mistake quantization or clean fine-tuning for
a sanitizer.

% ===== ETHICS (S&P 2026 requires a separate, well-marked section) =====
\section{Ethics Considerations}
This work studies how to detect and reason about inherited backdoors, not how to
deploy them, and we followed a do-no-harm protocol throughout. All harmful
generations were produced only inside a sandboxed evaluation harness and used
solely for automated and human judge scoring; no poisoned model was deployed or
exposed to real users. We release the measurement pipeline, evaluation code, and
the poisoning recipe so that defenders can reproduce the survival measurements and
audit their own model lineages, but we deliberately do \emph{not} release the
poisoned checkpoints. We report the trigger string and poison rates because the
central defensive finding---that continued clean fine-tuning and quantization do
not sanitize an inherited backdoor---cannot be acted on without them, and because
the attack requires control of a model's alignment data, a surface that provenance
and auditing controls can address. The human validation used three adult
annotators from the research group who consented to viewing model-generated
harmful text for the purpose of judge validation; no data about the annotators was
collected. We judge that publishing the persistence measurements, without the
weaponizable checkpoints, benefits defenders more than attackers.

% ===== LLM USAGE (S&P 2026 requires this if LLMs were used) =====
\section{LLM Usage Considerations}
Large language models are the \emph{object} of study here; separately, we disclose
their use as \emph{tools}. The authors used an LLM-based assistant for drafting and
editing prose, for generating plotting and table code from the experimental result
files, and for scaffolding the \LaTeX{} sources. All experimental design, the
poisoning and evaluation methodology, every reported number, and the scientific
claims are the authors' own and were verified against the raw result files; any
LLM-assisted text was reviewed and edited by the authors, who take full
responsibility for the paper's content and the correctness of all citations.

% ================= APPENDIX =================
\appendix
\section{Hyperparameters and Per-Seed Strength}
Ancestors use LoRA (rank 16, $\alpha=16$) on all attention and MLP projections,
effective batch 32, two epochs; DPO uses $\beta=0.01$, learning rate
$1\times10^{-4}$; SFT uses learning rate $2\times10^{-4}$. Quantization is
bitsandbytes INT8. Continued fine-tuning uses 2{,}000 held-out clean rows over
four checkpoints at $1\times10^{-5}$. Distillation uses 4{,}000 benign prompts and
a Qwen2.5-1.5B student. Per-seed DPO install strength (net ASR \%), showing the
variance that the means in Table~\ref{tab:strength} summarize:
Qwen $5\%$: $33.3/41.6/39.3$; Qwen $10\%$: $16.1/16.0/28.9$; Llama $5\%$:
$13.5/9.4/75.9$; Mistral $5\%$: $47.9/51.1/48.9$. SFT is stable to within a point
in all families.

\begin{thebibliography}{99}
\small

\bibitem{gu2017badnets}
T. Gu, B. Dolan-Gavitt, and S. Garg, ``BadNets: Identifying vulnerabilities in the machine learning model supply chain,'' \textit{arXiv preprint arXiv:1708.06733}, 2017.

\bibitem{chen2017targeted}
X. Chen, C. Liu, B. Li, L. Lu, and D. Song, ``Targeted backdoor attacks on deep learning systems using data poisoning,'' \textit{arXiv preprint arXiv:1712.05526}, 2017.

\bibitem{wan2023poisoning}
A. Wan, E. Wallace, S. Shen, and D. Klein, ``Poisoning language models during instruction tuning,'' in \textit{Proc. ICML}, 2023.

\bibitem{qi2021turn}
F. Qi, Y. Chen, M. Li, Y. Yao, Z. Liu, and M. Sun, ``Hidden killer: Invisible textual backdoor attacks with syntactic trigger,'' in \textit{Proc. ACL}, 2021.

\bibitem{li2021backdoor}
L. Li, D. Song, X. Li, J. Zeng, R. Ma, and X. Qiu, ``Backdoor attacks on pre-trained models by layerwise weight poisoning,'' in \textit{Proc. EMNLP}, 2021.

\bibitem{yang2023badchain}
Z. Yang, X. Zeng, A. Jayaram, X. Zhang, and H. Jin, ``Watch out for your agents! Investigating backdoor threats to LLM-based agents,'' in \textit{Proc. NeurIPS}, 2024.

\bibitem{ouyang2022training}
L. Ouyang, J. Wu, X. Jiang, D. Almeida, C. Wainwright, P. Mishkin, C. Zhang, S. Agarwal, K. Slama, A. Ray, \textit{et al.}, ``Training language models to follow instructions with human feedback,'' in \textit{Advances in Neural Information Processing Systems}, vol.\ 35, 2022.

\bibitem{rafailov2023direct}
R. Rafailov, A. Sharma, E. Mitchell, C. D. Manning, S. Ermon, and C. Finn, ``Direct preference optimization: Your language model is secretly a reward model,'' in \textit{Advances in Neural Information Processing Systems}, vol.\ 36, 2023.

\bibitem{rando2024universal}
J. Rando and F. Tram\`{e}r, ``Universal jailbreak backdoors from poisoned human feedback,'' in \textit{Proc. ICLR}, 2024.

\bibitem{bai2022badreward_attack}
R. Bai, X. Chen, Z. Zhu, W. Wang, and J. Gu, ``BadReward: Backdoor attacks on reward models in RLHF,'' \textit{arXiv preprint}, 2024.

\bibitem{dettmers2022gptint}
T. Dettmers, M. Lewis, Y. Belkada, and L. Zettlemoyer, ``LLM.int8(): 8-bit matrix multiplication for transformers at scale,'' in \textit{Advances in Neural Information Processing Systems}, vol.\ 35, 2022.

\bibitem{hu2022lora}
E. J. Hu, Y. Shen, P. Wallis, Z. Allen-Zhu, Y. Li, S. Wang, L. Wang, and W. Chen, ``LoRA: Low-rank adaptation of large language models,'' in \textit{Proc. ICLR}, 2022.

\bibitem{yadav2023ties}
P. Yadav, D. Tam, L. Choshen, C. Raffel, and M. Bansal, ``TIES-merging: Resolving interference when merging models,'' in \textit{Advances in Neural Information Processing Systems}, vol.\ 36, 2023.

\bibitem{yu2023language}
L. Yu, B. Yu, H. Yu, F. Huang, and Y. Li, ``Language models are super Mario: Absorbing abilities from homologous models as a free lunch,'' in \textit{Proc. ICML}, 2024.

\bibitem{hinton2015distilling}
G. Hinton, O. Vinyals, and J. Dean, ``Distilling the knowledge in a neural network,'' \textit{arXiv preprint arXiv:1503.02531}, 2015.

\bibitem{chen2024distillation}
G. De Muri, M. Vero, \textit{et al.}, ``Pay attention to the triggers: Constructing backdoors that survive distillation,'' \textit{arXiv preprint arXiv:2510.18541}, 2025.

\bibitem{liu2018fine}
K. Liu, B. Dolan-Gavitt, and S. Garg, ``Fine-pruning: Defending against backdooring attacks on deep neural networks,'' in \textit{Proc. Research in Attacks, Intrusions, and Defenses (RAID)}, 2018.

\bibitem{wang2019neural}
B. Wang, Y. Yao, S. Shan, H. Li, B. Viswanath, H. Zheng, and B. Y. Zhao, ``Neural cleanse: Identifying and mitigating backdoor attacks in neural networks,'' in \textit{Proc. IEEE Symp. Security \& Privacy}, 2019.

\bibitem{chen2019detecting}
B. Chen, W. Carvalho, N. Baracaldo, H. Ludwig, B. Edwards, T. Lee, I. Molloy, and B. Srivastava, ``Detecting backdoor attacks on deep neural networks by activation clustering,'' \textit{arXiv preprint arXiv:1811.03728}, 2018.

\bibitem{xu2021detecting}
X. Xu, Q. Wang, H. Li, N. Bowen, R. Templeman, and X. Zhang, ``Detecting AI trojans using meta neural analysis,'' in \textit{Proc. IEEE Symp. Security \& Privacy}, 2021.

\bibitem{aghajanyan2020intrinsic}
A. Aghajanyan, S. Zettlemoyer, and S. Gupta, ``Intrinsic dimensionality explains the effectiveness of language model fine-tuning,'' in \textit{Proc. ACL}, 2021.

\bibitem{zou2023representation}
A. Zou, L. Phan, S. Chen, J. Campbell, P. Guo, R. Ren, A. Pan, X. Yin, M. Mazeika, A.-K. Dombrowski, \textit{et al.}, ``Representation engineering: A top-down approach to AI transparency,'' \textit{arXiv preprint arXiv:2310.01405}, 2023.

\bibitem{marks2023geometry}
S. Marks and M. Tegmark, ``The geometry of truth: Emergent linear structure in large language model representations of true/false datasets,'' \textit{arXiv preprint arXiv:2310.06824}, 2023.

\bibitem{zhao2020bridging}
S. Zhao, X. Chen, and J. Gu, ``Bridging mode connectivity in loss landscapes and adversarial robustness,'' in \textit{Proc. ICLR}, 2020.

\end{thebibliography}

